\section{Comparación de herramientas de visualización}
\label{tools_comparison}
    \subsection{Panorama de herramientas}
    \label{tools_overview}
        Las herramientas de visualización pueden organizarse en cuatro enfoques,
        cada uno con un equilibrio distinto entre control, facilidad de uso,
        interactividad y reproducibilidad. El enfoque imperativo construye la 
        figura instrucción por instrucción y ofrece un control detallado sobre 
        cada elemento, en este grupo, Matplotlib proporciona gráficos de alta 
        resolución y una licencia libre, aunque requiere más código, de forma 
        similar, la graficación base de R funciona sin dependencias, pero ofrece 
        una composición por capas menos flexible \cite{castro2026ciencia}, 
        \cite{wilkinson2005grammar}, \cite{hunter2007matplotlib}. El enfoque 
        declarativo, inspirado en la gramática de gráficos de Wilkinson, 
        permite describir qué se desea mostrar y delega el trazado en la 
        herramienta, por ello, \textit{seaborn}, \textit{ggplot2} y 
        \textit{Altair} producen gráficos estadísticos expresivos con pocas 
        líneas de código, aunque con menor control sobre los detalles más 
        finos \cite{wilkinson2005grammar}. En el ámbito web, \textit{Plotly} 
        y \textit{D3.js} añaden interactividad de forma nativa, pero su desarrollo 
        suele requerir un esfuerzo mayor \cite{castro2026ciencia}. Finalmente, 
        las herramientas de interfaz gráfica reemplazan el código por la
        manipulación directa de los elementos, como es el caso de \textit{Tableau} 
        y \textit{Power BI} que permiten construir tableros rápidamente y 
        cuentan con una amplia adopción empresarial, sin embargo, ofrecen menor 
        reproducibilidad y requieren licencias comerciales \cite{castro2026ciencia}, 
        \cite{hunter2007matplotlib}.

    \subsection{Elección de herramientas: Matplotlib y R}
    \label{tools_choice}
        En Python se selecciona Matplotlib porque permite un mayor control sobre 
        los elementos de la figura, como la cuadrícula, las etiquetas y los ejes
        \cite{castro2026ciencia}, además, bibliotecas como \textit{seaborn} y el 
        método \texttt{DataFrame.plot} de \textit{pandas} utilizan Matplotlib 
        como base, por lo que su dominio facilita el uso de estas herramientas
        \cite{hunter2007matplotlib}. Su integración con \textit{pandas} y
        \textit{NumPy} también permite realizar los cálculos y generar las
        visualizaciones dentro de un mismo script reproducible, utilizando 
        funcionalidades de código abierto y sin costo \cite{mckinney2010data}, 
        \cite{harris2020array}. Por el lado de R, se utiliza la graficación base 
        como herramienta de contraste, pues repetir los mismos cálculos y 
        visualizaciones en un lenguaje diferente permite comprobar que los 
        principios de diseño no dependen de una herramienta específica y 
        ayuda a identificar posibles errores de implementación en cualquiera de 
        los dos entornos \cite{castro2026ciencia}.
